% Mẫu nhập nhiều lớp/chủ đề. Không cần dùng mã ID.
% Giữ cấu trúc ex, choice / choiceTF / shortans, loigiai.
% @grade áp dụng từ vị trí khai báo. @category có thể đặt riêng trong từng ex.
% @level: NB, TH, VD, VDC. Giáo viên kiểm tra lại mức độ trước khi lưu.
% Có thể dùng nhãn cũ %[10C4Ni-1T]; dạng nội dung lấy theo tiêu đề section.
% @grade: 10
\section*{Mệnh đề và phủ định}
\begin{ex}
% @level: NB
Mệnh đề nào dưới đây đúng?
\choice{\True $2+3=5$}{$2+3=6$}{$2+3=7$}{$2+3=8$}
\loigiai{Ta có $2+3=5$, nên chọn phương án A.}
\end{ex}

\section*{Các phép toán vectơ}
\begin{ex}
% @level: TH
Cho hình bình hành $ABCD$. Xét các khẳng định:
\choiceTF
{\True $\vec{AB}+\vec{AD}=\vec{AC}$}
{\True $\vec{AB}=\vec{DC}$}
{$\vec{AB}=\vec{CD}$}
{$\vec{AC}=\vec{0}$}
\loigiai{Hai ý đầu đúng theo quy tắc hình bình hành. Hai ý sau sai.}
\end{ex}

\begin{ex}
% @category: Tọa độ vectơ
% @level: TH
Cho $\vec{u}=(2;3)$. Tổng hai tọa độ của $\vec{u}$ bằng bao nhiêu?
\shortans{5}
\loigiai{Tổng hai tọa độ là $2+3=5$.}
\end{ex}
